(1) 取 \(y=e^{i(kx-\omega t)}\)，代入给 \(-\rho A\omega^2+EI k^4=0\)，所以 \(\omega=\sqrt{EI/(\rho A)}\,k^2\)，取正频率支。记 \(a_b=\sqrt{EI/(\rho A)}\)，其量纲为 \(\mathrm{m^2\,s^{-1}}\)；弯曲波是强色散波而非恒速拉伸弦波。
(2) 正波数时相速度 \(v_p=\omega/k=a_b k\)，群速度 \(v_g=d\omega/dk=2a_bk=2v_p\)，二者都随 \(k\) 线性增长。对于负波数，群速度带传播方向符号 \(v_g=2a_bk\)；在高频下欧拉梁模型忽略剪切与转动惯性会失效，不可把群速度无限增长当真实物理极限。
(3) 取简支空间本征函数 \(Y_n(x)=\sin(n\pi x/L)\)，其位移及二阶导在两端均为零，\(n=1,2,\ldots\)。代入方程给 \(\omega_n=a_b(n\pi/L)^2\)，因此频率比 \(\omega_n/\omega_1=n^2\)，而不是张紧弦的整数倍 \(n\)。
(4) 总能量 \(H=\frac12\int_0^L[\rho A y_t^2+EI y_{xx}^2]\,dx\)。对时间求导，利用方程与两次分部积分，边界功率项为 \([EI(y_{xx}y_{xt}-y_{xxx}y_t)]_0^L=0\)，因为简支端 \(y_t=0\) 且 \(y_{xx}=0\)；故 \(\dot H=0\)。最低频率 \(\omega_1\propto L^{-2}\)，加长梁显著降低固有振动频率。
